\documentclass[journal=jpclcd, manuscript=letter, layout=twocolumn, doi=true]{achemso}

\usepackage[utf8]{inputenc}
\usepackage{amsmath}
\usepackage{amssymb}
\usepackage{graphicx}
\usepackage{physics}
\usepackage{booktabs}
\usepackage{xcolor}
\usepackage[colorlinks=true, allcolors=blue]{hyperref}
\usepackage{subcaption}

\newcommand{\red}{\color[rgb]{0.8,0,0}}
\newcommand{\green}{\color[rgb]{0.0,0.6,0.0}}
\newcommand{\blue}{\color[rgb]{0.0,0.0,0.6}}

\title{Non-Equilibrium Dynamics of the Time-Dependent Excitonic Coupling in Protein Dimers}

%\title{Interplay of Time-Dependent Excitonic Coupling and Decoherence in Protein Dimers}

\author{Robson Christie}
\affiliation{School of Mathematics and Physics, University of Portsmouth, PO1 3FX, United Kingdom}

\author{Cerys Murray}
\affiliation{Leverhulme Quantum Biology Doctoral Training Centre, School of Biosciences, Advanced Technology Institute, University of Surrey, Guildford GU2 7XH, UK}

\author{Youngchan Kim}
%\email{youngchan.kim@surrey.ac.uk}
\affiliation{Leverhulme Quantum Biology Doctoral Training Centre, School of Biosciences, Advanced Technology Institute, University of Surrey, Guildford GU2 7XH, UK}

\author{Jaewoo Joo}
\email{* jaewoo.joo@port.ac.uk}
\affiliation{School of Mathematics and Physics, University of Portsmouth, PO1 3FX, United Kingdom}

\begin{document}

\maketitle

% Move tocentry here to bypass layout-induced frame distortion
\begin{tocentry}
\includegraphics[width=1.0\textwidth]{Structure3.pdf}
\end{tocentry}

\begin{abstract}
We quantify the excitonic coupling within the Venus\textsubscript{A206} protein dimer using a quantum-classical hybrid workflow. By employing a transition-density coupling formalism, we calculate a coupling strength of $J = 74.38~\mathrm{cm^{-1}}$, which is approximately 5.6 times stronger than the far-field point-dipole estimate of $13.31~\mathrm{cm^{-1}}$. This disparity highlights the role of near-field multipolar effects at the 27.6~\AA\ chromophore centroid separation. Furthermore, we analyse the impact of non-equilibrium dielectric screening, demonstrating that the optical dielectric constant ($\varepsilon_{\rm opt}$), rather than the static limit ($\varepsilon_{\rm eq}$), governs the initial excitonic interaction. This non-equilibrium approach provides a theoretical framework for the robust Davydov splitting observed in circular dichroism spectra. This spectroscopic signature is preserved because the excitation wavelength ($\sim 515$~nm) is orders of magnitude larger than the dimer separation, leading to a collective excitation where the splitting is imprinted at the moment of photon absorption, preceding the rapid environmental dephasing that occurs on sub-picosecond timescales. To study this post-absorption evolution, we employ Markovian stochastic dynamics to simulate the transition from the initial exciton superposition over both monomers toward a localised state, followed by subsequent incoherent hopping between these localised states. Future work will extend this framework to incorporate explicit vibronic coupling and evaluate the impact of non-Markovian protein phonons on long-term coherence dynamics.
\end{abstract}

\section{Introduction}

Excitation energy transfer in biological systems is conventionally understood through the competition between intermolecular excitonic coupling ($J$) and environmental vibrational dephasing \cite{forster1965delocalized, may2004charge}. When environmental fluctuations dominate, phase relationships randomise rapidly, restricting the system to the weak-coupling regime characterised by incoherent Förster hopping \cite{forster1948intermolecular}. Conversely, the intermediate and strong regimes permit the formation of coherently delocalised excitonic eigenstates, manifesting spectroscopically as Davydov splitting \cite{davydov1971}. Because physiological environments are highly dynamic and structurally disordered, biophysics has historically assumed that ambient thermal noise precludes robust excitonic delocalisation, confining such systems to the weak interaction regime \cite{kim2021quantum}.

Homodimers of the yellow fluorescent protein variant Venus\textsubscript{A206} challenge this assumption. Operating at physiological temperatures, these dimers exhibit distinct Davydov splitting in circular dichroism (CD) spectra, and photon antibunching statistics confirm that the interacting chromophores act as a collective quantum entity upon photon absorption \cite{kim2019venusa206, freed2025effect}. These signatures indicate an interaction strength firmly within the intermediate excitonic regime, seemingly at odds with the biological environment inducing strong decoherence.

Reconciling the enhanced excitonic coupling with strong environmental noise requires separating the physical timescales of exciton formation and subsequent dephasing. Because the excitation wavelength ($\sim 515$~nm) exceeds the inter-chromophore separation ($\sim 2.74$~nm) by two orders of magnitude, photoexcitation acts as a collective, near-instantaneous vertical transition directly into a delocalised bright eigenstate. The Davydov splitting is an inherent property of the energy spectrum imprinted at this precise moment of absorption. As demonstrated by the stochastic trajectories in this work, the system subsequently undergoes strong environmental dephasing and rapid relaxation into a localised state. However, because this decoherence and the associated dielectric screening occur on sub-picosecond to picosecond timescales, they are significantly slower than the absorption event. Consequently, strong excitonic coupling and its spectroscopic signatures are observed despite the presence of strong, subsequent decoherence.

Accurately predicting the magnitude of the coupling $J$ that drives this initial delocalisation requires overcoming the limitations of standard point dipole interaction models. At the chromophore separation in Venus\textsubscript{A206}, the point-dipole approximation severely underestimates the interaction. At these intermolecular distances, near-field multipolar effects become dominant contributors to the overall coupling strength \cite{scholes2003, kitohnishioka2017forster}. Capturing these effects requires integrating the full three-dimensional spatial distribution of the electron transition density, which maps the shape of the transition's probability amplitude onto the molecular structure. 

Furthermore, because exciton formation precedes the macroscopic nuclear reorganisation of the solvent (Debye relaxation $\tau_D \approx 8.3~\mathrm{ps}$), the initial Coulombic coupling is not damped by the static permittivity of water ($\varepsilon_{\mathrm{eq}} \approx 78$). Instead, it is screened only by the fast electronic polarisation of the environment, corresponding to the optical dielectric limit ($\varepsilon_{\mathrm{opt}} \approx 1.78$) \cite{mennucci2012, scholes2003}.

In this work, we present a quantum-classical hybrid analysis of the Venus\textsubscript{A206} dimer. By combining a transition-density coupling (TDC) formalism with a non-equilibrium dielectric framework, we quantify the massive near-field multipolar contributions to the excitonic coupling. We demonstrate that the interplay of optical-limit screening and transition-density geometry is sufficient to produce the experimentally observed intermediate coupling, independent of the subsequent incoherent hopping dynamics.

%%%%%%%%%%%%%%%%%%%%%%%%%%%%%%%%%%%%%%%%%%%%%%%%
\section{Theoretical Framework}

\begin{figure*}[t]
  \centering
  \begin{subfigure}[b]{0.32\textwidth}
    \centering
    \includegraphics[width=\linewidth]{Fig_Bloch_Grid.pdf}
    \caption{}
    \label{fig:bloch}
  \end{subfigure}\hfill
  \begin{subfigure}[b]{0.32\textwidth}
    \centering
    \includegraphics[width=\linewidth]{Fig_SSE_Site.pdf}
    \caption{}
    \label{fig:sse_site}
  \end{subfigure}\hfill
  \begin{subfigure}[b]{0.32\textwidth}
    \centering
    \includegraphics[width=\linewidth]{Fig_SSE_Adiabatic.pdf}
    \caption{}
    \label{fig:sse_adiabatic}
  \end{subfigure}

  \caption{Simulation of excitonic dynamics in the Venus dimer with the local site states $\ket{1}\equiv\ket{e_1 g_2}$ and $\ket{2}\equiv\ket{g_1 e_2}$: (a) Unified Bloch sphere showing the stochastic pure-state trajectory (red) and ensemble-averaged mixed-state path (blue); (b) Site-basis populations and coherence for a single SSE trajectory; (c) Adiabatic populations of bright and dark states $\ket{\pm}$ for a single SSE trajectory.}
  \label{fig:ExcitonDynamics}
\end{figure*}

To describe the excitonic interactions within the dimer, we restrict our model to the single-exciton manifold, defining the local site basis as $\ket{1}\equiv\ket{e_1 g_2}$ and $\ket{2}\equiv\ket{g_1 e_2}$. This notation explicitly enforces the conservation of energy in a single-photon absorption subspace, ensuring that only a single excitation is shared between the two chromophores. Because the $\sim 15$~\AA\ minimum separation between the chromophores precludes significant short-range Dexter exchange interactions~\cite{skourtis2016dexter}. Consequently, the electronic coupling $J$ is accurately described as the Coulomb interaction between the transition densities of the individual monomers \cite{hsu2009}. The electron transition density, $\rho^{g\to e}_m(\mathbf r)$, is obtained by integrating the product of the ground and excited state wavefunctions over the coordinates of $N-1$ electrons in a monomer:
\begin{multline} \label{eq:transition_density}
    \rho^{g\to e}_m(\mathbf{r}) =  
    N \int \dots \int d\mathbf{r}_2 \dots d\mathbf{r}_N \\ \Psi_e^*(\mathbf{r}, \mathbf{r}_2, \dots, \mathbf{r}_N) \Psi_g(\mathbf{r}, \mathbf{r}_2, \dots, \mathbf{r}_N) \,
\end{multline}
where $\Psi_g$ and $\Psi_e$ are the many-body electronic wavefunctions for the ground and excited states, respectively. Crucially, while the standard electron density is a positive-definite observable representing the probability of finding an electron, the transition density is an off-diagonal matrix element.
Crucially, while the standard electron density is a positive-definite observable representing the probability of finding an electron, the transition density is an off-diagonal matrix element. It integrates to zero over all space ($\int \rho^{g\to e}(\mathbf{r}) d\mathbf{r} = 0$) due to the orthogonality of the ground and excited states, and its spatial distribution, characterised by positive and negative phases, defines the transition's multipolar character. Although this zero-th moment evaluates to zero, weighting the transition density by the position vector yields the transition dipole moment ($\boldsymbol{\mu} = \int \mathbf{r} \rho^{g \to e}(\mathbf{r}) d\mathbf{r}$), which dictates the far-field optical response. The full TDC is given by the spatial integral
\begin{equation}
    J_{\mathrm{TDC}} \approx \frac{1}{4\pi\varepsilon(t)}\iint \frac{\left(\rho^{g\to e}_1(\mathbf r)\right)^*\,\rho^{g\to e}_2(\mathbf r')}{|\mathbf r-\mathbf r'|} \,d\mathbf r\,d\mathbf r' \,
    \label{eq:J_integral}
\end{equation}
for time-dependent permittivity $\varepsilon(t)$.
Formally, the transition density can be complex, and the conjugate ensures a real coupling strength. In our implementation, the space is discretized over a grid as $J \approx\frac{1}{4\pi\varepsilon(t)}\sum_{i\in 1}\sum_{j\in 2}\frac{q_i q_j}{|\mathbf r_i-\mathbf r_j|}$, where $q_i$ and $q_j$ represent the discrete transition atomic charges. For comparison, we also evaluate the far-field point-dipole approximation (PDA), which collapses these densities into two idealized dipoles $\boldsymbol{\mu}_m$ separated by vector $\mathbf{R}$ given by
\begin{equation} \label{eq:J_PD}
    J_{\mathrm{PDA}} = \frac{1}{4\pi\varepsilon(t)} \left( \frac{\boldsymbol{\mu}_1 \cdot \boldsymbol{\mu}_2}{|\mathbf{R}|^3} - \frac{3(\boldsymbol{\mu}_1 \cdot \mathbf{R})(\boldsymbol{\mu}_2 \cdot \mathbf{R})}{|\mathbf{R}|^5} \right) .
\end{equation}

A critical factor is the time-dependent effective permittivity, $\varepsilon(t)$. The aqueous environment exhibits a response characterised by a Debye relaxation time $\tau_D \approx 8.3~\mathrm{ps}$. Because this is orders of magnitude longer than the sub-picosecond timescales of excitation transfer, the solvent cannot undergo nuclear equilibration. Thus, the time-dependent permittivity is initially governed by the optical dielectric constant ($\varepsilon_{\rm opt} \approx 1.78$), and the screening effect approaches the static limit ($\varepsilon_{\rm eq} \approx 78$) as given by
\begin{equation}
    \frac{1}{\varepsilon(t)} = \frac{1}{\varepsilon_{\rm eq}} + \left(\frac{1}{\varepsilon_{\rm opt}} - \frac{1}{\varepsilon_{\rm eq}}\right)e^{-t/\tau_D}.
\end{equation}

Assuming a structurally identical homodimer the ground to first excited state excitation energies are identical for each monomer ($E_1 = E_2 = E$), the excited state Hamiltonian in the site basis is
\begin{equation}
    \hat H_{\text{site}}(t)=\begin{pmatrix} E & J(t) \\ J(t) & E \end{pmatrix}.
\end{equation}
Diagonalizing $\hat H_{\text{site}}(t)$ yields the adiabatic (exciton) basis, where eigenstates are the symmetric (bright) and antisymmetric (dark) combinations: $\ket{\pm} = \frac{1}{\sqrt{2}}(\ket{1} \pm \ket{2})$. Photoexcitation is treated as a vertical transition directly into the delocalised bright state $\ket{+}$. The energy separation is the Davydov splitting, $\Delta E(t) = 2|J(t)|$. Over several picoseconds, the coupling $J(t)$ relaxes as the solvent environment reorganises from the initial optical dielectric response toward the static equilibrium limit, resulting in a gradual damping of the excitonic interaction.

Because the physical interactions between the system and its aqueous environment are local, we define a Hermitian jump operator $\hat{L} = \sqrt{\hbar \gamma / 2} \hat{\sigma}_z$ in the site basis, where $\gamma = 1/T_2^*$ is the pure dephasing rate ($\hat{\sigma}_z$: Pauli-$z$ operator). For our simulations, a dephasing time of $T_2^* = 60~\mathrm{fs}$ was used ($\gamma \approx 16.67~\mathrm{ps^{-1}}$) to reflect the rapid environmental decoherence. This operator effectively ``measures'' the excitation site, driving the state towards the local site basis, which emerges as the decoherence resistant pointer basis. The system's ensemble-averaged density matrix $\hat{\rho}(t)$ evolves via the Lindblad equation \cite{manzano2020short}
\begin{equation}
    \frac{d\hat{\rho}}{dt}=-\frac{i}{\hbar}[\hat H(t),\hat\rho] + \hat{L} \hat\rho \hat{L} - \frac{1}{2} \{\hat{L}^2, \hat\rho\}
    \label{eq:ME01}
\end{equation}

Building on the Lindblad model of \cite{abrahams2025excitonic}, we extend the analysis by simulating individual pure-state trajectories in addition to the ensemble density operator. While the master equation (ME) provides a deterministic description of the average behaviour of a large ensemble of dimers, it cannot resolve the discrete, incoherent events experienced by a single system. To capture these dynamics, we unravel the ME into a stochastic Schrödinger equation (SSE), which describes the evolution of a single pure-state trajectory $|\psi\rangle$ subject to continuous environmental monitoring. The density operator is recovered by taking the stochastic average over many such trajectories, $\hat{\rho} = \mathbb{E}[|\psi\rangle\langle\psi|]$. The SSE is given by
\begin{multline}
    |d\psi\rangle = \left( -\frac{i}{\hbar}\hat{H} - \frac{1}{2}(\hat{L} - \langle \hat{L} \rangle)^2 \right) |\psi\rangle dt \\
    + (\hat{L} - \langle \hat{L} \rangle) |\psi\rangle dW
\end{multline}
where $dW$ is an It\^{o} Wiener process~\cite{petruccione}.

We visualise the individual pure-state trajectories (\textbf{Figures~\ref{fig:ExcitonDynamics}(b) and (c)}) by mapping the system to a Bloch vector $\mathbf{r} = (u, v, w)$ (\textbf{Figure~\ref{fig:ExcitonDynamics}(a)}), where $u = 2\text{Re}(\rho_{12})$, $v = 2\text{Im}(\rho_{12})$, and $w = \rho_{11} - \rho_{22}$. All the numerical simulations are initialised in the delocalised bright state $\ket{+}$, corresponding to a vector on the equator of the Bloch sphere. Individual trajectories exhibit incoherent hopping between sites 1 and 2; in the bright-dark basis, this manifests as an alternation between the bright state and superpositions of bright and dark states. These stochastic paths represent pure-state evolution restricted to the surface of the sphere. Conversely, the ensemble-averaged dynamics (governed by the ME in Eq.~(\ref{eq:ME01})) describe the rapid, smooth decay of off-diagonal coherence, leading to a statistical mixture over both sites and evolution toward the centre of the Bloch sphere.

\section{TDDFT Results}

\begin{table}[h]
\centering
\caption{TDDFT excited states for a monomer used in the coupling analysis. The wavelengths indicate the photon wavelengths corresponding to TDDFT transition energies $E$ computed in Terachem.}
\label{tab:roots}
\begin{tabular}{rccc}
\toprule
State & Wavelength (nm) & $f$ & $|\mu|$ (a.u.) \\
\midrule
1  & 542.4 & 0.7037 & 3.5449 \\
2  & 476.9 & 0.0238 & 0.6113 \\
3  & 460.2 & 0.0030 & 0.2132 \\
4  & 443.9 & 0.0448 & 0.8091 \\
5  & 391.3 & 0.0304 & 0.6258 \\
6  & 388.8 & 0.0029 & 0.1927 \\
7  & 380.9 & 0.0310 & 0.6235 \\
8  & 340.0 & 0.0304 & 0.5833 \\
9  & 334.7 & 0.0000 & 0.0000 \\
10 & 329.9 & 0.0033 & 0.1893 \\
\bottomrule
\end{tabular}
\end{table}

Ten Time-Dependent Density Functional Theory (TDDFT) roots (as eigenenergies) were computed using TeraChem, as shown in Table \ref{tab:roots}. Root 1 was identified as the target excitonic state, located at $542.4~\mathrm{nm}$ with a high oscillator strength of $f=0.7037$ and a transition dipole moment of $|\mu|=3.5449~\mathrm{a.u.}$. To physically characterise this state, the transition density and the ground-to-excited state electron density difference were mapped onto the dimer frame (\textbf{Figure~\ref{fig:density_combined}} and \textbf{Figure~\ref{fig:electron_combined}}, respectively). 

The electron transition density in \textbf{Figure~\ref{fig:density_combined}}, which serves as the fundamental quantity for evaluating the Coulombic integral in Eq.~(\ref{eq:J_integral}), exhibits a complex multipolar spatial distribution that extends across the entire chromophore and adjacent residues. Note that the red and blue isosurfaces in \textbf{Figure~\ref{fig:density_combined}(b)} do not represent static charge gain or loss, but rather the spatial ``lobes'' of the transition probability. The spatial overlap of these lobes between the two chromophores determines the magnitude of $J_{\text{TDC}}$, capturing interactions that are ignored when the density is collapsed into a single point-dipole vector. Concurrently, the electron density difference in \textbf{Figure~\ref{fig:electron_combined}(b)} illustrates the charge reorganisation from the ground to the first excited state, with the deprotonated phenolic oxygen on the chromophore residue (CR2, anionic form) transferring charge to the rest of the chromophore. 

The transition densities were mapped to the dimer frame (centroid separation $27.6~\mathrm{\AA}$). The inter-dipole angle was calculated as $92.85^\circ$. The $t=0$ optical dielectric coupling results are summarised below:

\begin{itemize}
    \item \textbf{Transition Density Coupling:} \\ $J_{\text{TDC}}(0) = 74.38~\mathrm{cm^{-1}}$ ($9.22~\mathrm{meV}$).
    \item \textbf{Point-Dipole Estimate:} \\ $J_{\text{PDA}}(0) = 13.31~\mathrm{cm^{-1}}$ ($1.65~\mathrm{meV}$).
\end{itemize}

\begin{figure*}[t]
  \centering
  \begin{subfigure}[b]{0.48\textwidth}
    \centering
    \includegraphics[width=\linewidth]{dimer_view_old_Difference_current.png}
    \caption{}
    \label{fig:density_dimer}
  \end{subfigure}\hfill
  \begin{subfigure}[b]{0.48\textwidth}
    \centering
    \includegraphics[width=\linewidth]{QMRegion.png}
    \caption{}
    \label{fig:density_qm}
  \end{subfigure}
  \caption{Visualization of the electron transition density in Eq.~(\ref{eq:transition_density})  between the ground and first excited state in Table \ref{tab:roots}. The red and blue isosurfaces represent the positive and negative phases of the transition density, mapped onto the dimer frame using PyMOL structural alignment. (a) The full dimer system partitioned into QM and MM domains. The MM embedding environment, treated as point charges, is rendered as a protein backbone with a rainbow colourmap to give a sense of depth perception. (b) A detailed view of the QM region (sticks and spheres) on the right monomer. This subset is defined by the chromophore, explicit nearest-neighbour residues, and water molecules, allowing for the treatment of electronic polarisation and deprotonation effects. PyMOL \cite{delano2002pymol} generation scripts are available on GitHub.}
  \label{fig:density_combined}
\end{figure*}
\begin{figure*}[t]
  \centering
  \begin{subfigure}[b]{0.48\textwidth}
    \centering
    \includegraphics[width=\linewidth]{dimer_view_old_Electron_current.png}
    \caption{}
    \label{fig:electron_dimer}
  \end{subfigure}\hfill
  \begin{subfigure}[b]{0.48\textwidth}
    \centering
    \includegraphics[width=\linewidth]{QMRegionElectron.png}
    \caption{}
    \label{fig:electron_qm}
  \end{subfigure}
  \caption{Visualization of the electron density difference between the ground and first excited state delocalised over the Venus\textsubscript{A206} dimer geometry from Table \ref{tab:roots}. The green and violet isosurfaces represent electron gain and loss respectively. (a) The dimer frame illustrating the QM/MM partitioning, with the MM backbone rendered in the rainbow colourmap to give a sense of depth perception. (b) A detailed subset of the QM region on the right monomer, highlighting the chromophore and explicit nearest-neighbour residues. The detailed PyMOL \cite{delano2002pymol} codes are available on GitHub.}
  \label{fig:electron_combined}
\end{figure*}

The computed coupling strength of $J=74.38~\mathrm{cm^{-1}}$, which yields a theoretical Davydov splitting of $\Delta E \approx 149~\mathrm{cm^{-1}}$, is in better qualitative agreement with recent experimental measurements. Specifically, Nguyen et al.\ reported an apparent coupling strength between $131~\mathrm{cm^{-1}}$ and $186~\mathrm{cm^{-1}}$ derived from the CD spectra of a dVenus tandem dimer \cite{nguyen2025anomalous}. Both our full TDC calculations and the experimental CD observations firmly place the Venus dimer interaction within the intermediate excitonic coupling regime ($10 - 1000~\mathrm{cm^{-1}}$) \cite{nguyen2025anomalous, kim2019venusa206}. Further improvements to the simulation estimates are theoretically possible, though they would incur a significant computational cost. Such refinements include enlarging the QM region and fully QM treating phenolic hydrogen to capture the tunnelling effects of protons \cite{chakraborty2008development,bourne2024quantum}.

Furthermore, the disparity between the full TDC calculation ($74.38~\mathrm{cm^{-1}}$) and the far-field point-dipole estimate ($13.31~\mathrm{cm^{-1}}$) provides a rigorous QM rationale for the robust nature of this coupling. Our structural model's centroid separation of $27.6~\mathrm{\AA}$ aligns closely with the $\sim 25~\mathrm{\AA}$ centre-to-centre chromophore proximity measured in Venus\textsubscript{A206} crystals \cite{kim2019venusa206}. At these distances, the point-dipole approximation severely underestimates the interaction by a factor of $\sim 5.6$. While the classical point-dipole model predicts a rapid inverse-cube distance decay, the spatial extent of the transition density in \textbf{Figure~\ref{fig:density_combined}} allows specific phase lobes to interact at effective distances much shorter than the centroid separation.

\section{Computational Method}

The computational workflow couples OpenMM \cite{openmm2017} for classical structural and solvent relaxation and TeraChem \cite{terachem2009,seritan2021terachem} for excited-state electronic structure calculations. The starting monomer structure (PDB: 1MYW) was hydrogenated and solvated with a 10.0~\AA\ padding of water molecules using PDBFixer and OpenMM. The system was parameterised using the AMBER14~\cite{amber14_manual} force field. An energy minimisation of the geometry was performed in OpenMM to relax steric clashes.

The QM region was constructed around the CR2 residue. The QM selection criteria included the full CR2 residue and the five nearest-neighbour residues and five nearest water molecules. The boundary between the QM and MM regions was treated using a link-atom scheme with hydrogen caps placed along the covalent bond vectors at the QM/MM interface. The final QM region comprised 238 atoms. To avoid unphysical polarisation of the QM electron density by proximal point charges, MM charges within a 1.8--2.8~\AA\ buffer zone of the QM region were smoothly damped. A charge-redistribution algorithm was applied to the affected residue groups (52 total) to ensure that the net charge of the MM embedding remained constant ($|\Delta Q|_{\mathrm{total}}=0$), where $\Delta Q$ represents the change in partial charge of the modified MM atoms, while preserving local residue neutrality.

TDDFT calculations were performed on the classically relaxed monomer geometry. We employed the range-separated hybrid functional $\omega$B97X-D3 with the 6-311G** basis set \cite{mardirossian2017}. The entire pipeline was executed in approximately two hours on a single NVIDIA RTX 4080 GPU, with the TDDFT computations dominating the total processing time. The effects of the bulk solvent on the electronic ground and excited states were treated using the COSMO Polarisable Continuum Model (PCM) \cite{seritan2021terachem} with microsolvation incorporated by including the five nearest water molecules to the CR2 residue in the QM region.

Ten monomer excited states were computed, and the brightest state in the 400--700~nm visible window was selected for analysis. The dimer model was generated from the PDB monomer by applying the crystallographic symmetry and scaling operations embedded in the 1MYW coordinate file to the relaxed structure and transition density. 

The excitonic coupling $J$ was calculated via Eq.~\eqref{eq:J_integral} over the dimer. While the electronic structure was solved in the presence of a high-dielectric continuum ($\varepsilon_{\rm eq}\approx 78$) to represent the equilibrated ground state, the Coulombic coupling between monomers was screened by the optical dielectric constant ($\varepsilon_{\rm opt} \approx 1.78$), consistent with the fast timescale of excitonic transfer relative to solvent nuclear reorganisation.

\section{Summary}

We have characterised the excitonic coupling and energy transfer dynamics in a Venus\textsubscript{A206} dimer using a hybrid QM/MM transition-density approach and open quantum systems modelling. Electronic structure calculations predict a robust coupling strength of $J = 74.38~\mathrm{cm^{-1}}$, corresponding to a Davydov splitting of $\Delta E \approx 149~\mathrm{cm^{-1}}$. By integrating distributed transition densities, we demonstrate that the standard point-dipole approximation underestimates the interaction by a factor of $\sim 5.6$, necessitating the explicit treatment of near-field multipolar effects at biologically relevant distances.

Our analysis validates the role of non-equilibrium dielectric screening in protein environments. Because the excitonic interaction timescale is significantly faster than the macroscopic Debye relaxation time of bulk water ($\sim 8.3~\mathrm{ps}$), the coupling operates in the optical dielectric limit ($\varepsilon_{\rm opt}$). This inertial ``freezing'' of the solvent protects the initial coupling from overwhelming equilibrium screening ($\varepsilon_{\rm eq}$). The stochastic trajectories illustrate that this coherence is transient; environmental interactions drive the system into a localised pointer basis on a sub-$100$~fs timescale. This temporal separation reconciles the robust Davydov splitting imprinted at absorption with rapid localization and subsequent screening within picosecond timescales.

\section{Data and Code Availability}

The computational workflow was executed using the script \texttt{terachem\_full\_pipeline.py} available on GitHub: \url{https://github.com/rchristie95/TerachemVenusA206}. 

\section{Acknowledgements}

\bibliography{Bibliography}

\end{document}